\section{目录}

\begin{itemize}
\item 1 数据一致性模型
  \begin{itemize}
  \item 1.1 顺序一致性
  \item 1.2 线性一致性
  \item 1.3 因果一致性
  \item 1.4 最终一致性
  \end{itemize}
\item 2 客户端一致性模型
\item 3 副本管理 (对应 7.4 Replica management)
\end{itemize}

复制 (Replication) 为什么复制？
\begin{itemize}
\item 提高可靠性：如果一个副本失效，可以切换到另一个副本。
\item 提高性能与可扩展性：分散请求负载，或利用地理位置接近性实现快速响应。
\end{itemize}

现实问题：
\begin{itemize}
\item 拥有多个副本意味着当任何副本发生更改时，该更改必须在所有副本上执行。
\item 副本需要保持一致。
\end{itemize}

性能与可扩展性主要问题：
\begin{itemize}
\item 为了保持副本一致性，通常需要确保所有冲突操作在各处以相同顺序执行。
\item 冲突操作：来自事务领域的概念。
  \begin{itemize}
  \item 读-写冲突：读操作和写操作并发执行。
  \item 写-写冲突：两个并发写操作。
  \end{itemize}
\end{itemize}

挑战：保证不发生冲突操作的全局排序可能代价高昂，会降低可扩展性。
解决方案：弱化一致性要求，以避免全局同步。

\section{数据一致性模型}

数据一致性模型
一致性模型：
\begin{itemize}
\item 分布式数据存储与进程之间的契约。
\item 精确说明在并发情况下读写操作的结果。
\end{itemize}

\begin{verbatim}
Process   Process   Process
  |         |         |
Local copy  Local copy  Local copy
  \         |         /
   Distributed data store
\end{verbatim}

符号约定
读写操作符号：
\begin{itemize}
\item $W_i(x)a$：进程 $P_i$ 向数据项 $x$ 写入值 $a$。
\item $R_i(x)b$：进程 $P_i$ 从数据项 $x$ 读出值 $b$。
\end{itemize}
所有数据项初始值为 NIL。

示例：$P_i: W(x)a \;\rightarrow\; P_j: R(x)\text{NIL} \;\rightarrow\; R(x)a$

\subsection{顺序一致性}

顺序一致性 (Sequential Consistency) 定义：
\begin{itemize}
\item 所有结果需等同于所有进程按某种特定顺序执行，且
\item 该顺序执行结果要和每个进程自己代码里写的操作顺序一致。
\end{itemize}

例子：一个共享的在线文档，虽然每个人的更新到达时间可能不同，但系统最终会保证所有人看到的修改历史是同一个顺序 (比如：A 先改，然后 B 改，然后 C 改)，并且 A 的修改一定发生在 B 之前，B 在 C 之前。

顺序一致性示例。三个并发进程 (初始值: 0)：
\begin{itemize}
\item P1: x = x + 1; print(y, z);
\item P2: y = y + 1; print(x, z);
\item P3: z = z + 1; print(x, y);
\end{itemize}

执行序列：

\begin{verbatim}
Execution 1   Execution 2   Execution 3   Execution 4
P1: x=1;      P1: x=1;      P2: y=1;      P3: z=1;
P2: y=1;      P3: z=1;      P1: x=1;      P1: x=1;
P3: z=1;      P2: y=1;      P3: z=1;      P2: y=1;
print(x,y,z); print(x,y,z); print(x,y,z); print(x,y,z);
Prints: 001011  Prints: 101011  Prints: 010111  Prints: 111111
Signature:001011 | Signature:101011 | Signature:110101 | Signature:111111
 (a)             (b)              (c)              (d)
\end{verbatim}

小测验：看如下 P1 与 P2 并发的例子：这个运行结果满足顺序一致性吗？不满足，因为我们无法找到某个能产生这样效果的全局序。

\subsection{线性一致性}

线性一致性 (Linearizability) 定义：
\begin{itemize}
\item 每个操作都应该看起来在其开始和完成之间的某个时刻瞬间生效。
\end{itemize}

与顺序一致性的区别：
\begin{itemize}
\item 线性一致性要求操作具有实时性 (原子性)，即操作的效果看起来是瞬间发生的。
\item 顺序一致性只要求操作的最终顺序，不要求实时性。
\item 线性一致性 = 顺序一致性 + 实时性约束。要求每个操作在调用与返回之间的某个“瞬间”原子生效，且跨进程的操作顺序必须与真实时间先后严格一致。
\item 顺序一致性只要求：所有节点看到相同的全局操作总序；该顺序不违反单个进程内的程序执行顺序。不要求符合跨进程的真实时间先后。
\end{itemize}

示例对比：

\begin{table}[htbp]
\centering
\caption{线性一致性示例}
\begin{tabular}{ll}
\toprule
Result & Possible ordering of operations \\
\midrule
R1(x)b | R2(y)b & W1(x)a; W1(y)a; W2(y)b; W2(x)b \\
R1(y)b | R2(x)a & W1(x)a; W2(y)b; W1(y)a; W2(x)b \\
R1(x)b | R2(y)a & W1(x)a; W2(y)b; W2(x)b; W1(y)a \\
R1(y)b | R2(x)a & W2(y)b; W1(x)a; W1(y)a; W2(x)b \\
R1(x)b | R2(y)a & W2(y)b; W1(x)a; W2(x)b; W1(y)a \\
R1(x)a | R2(y)a & W2(y)b; W2(x)b; W1(x)a; W1(y)a \\
\bottomrule
\end{tabular}
\end{table}

\subsection{因果一致性}

因果一致性 (Causal Consistency)：
\begin{itemize}
\item 潜在因果相关的写操作必须被所有进程以相同顺序看到。
\item 并发的写操作可能被不同进程以不同顺序看到。
\end{itemize}

(a) 因果不一致的数据存储 (因为 P2 的 R(x)a 对 W(c)b 产生影响，所以 W(a)a -> W(c)b)
\begin{verbatim}
P1: W(x)a
P2: R(x)a  W(x)b
P3: R(x)a  R(x)b
P4:        R(x)a R(x)b
\end{verbatim}
(b) 因果一致的数据存储
\begin{verbatim}
P1: W(x)b
P2: R(x)b R(x)a W(x)b
P3: R(x)a R(x)b
P4: R(x)a R(x)b
\end{verbatim}

可串行化 / 事务可串行化 (serializability)：对一组事务中的所有操作进行排序，使得最终结果与这些事务的某种串行执行相匹配。

示例：
\begin{verbatim}
BEGIN_TRANSACTION
x=0
x=x+1
END_TRANSACTION

BEGIN_TRANSACTION
x=0
x=x+2
END_TRANSACTION

BEGIN_TRANSACTION
x=0
x=x+3
END_TRANSACTION

Transaction T1  Transaction T2  Transaction T3
A number of schedules
Time ->
S1 | x=0
S2 | x=0
S3 | x=0
S4 | x=0
x=x+1
x=0
x=0
x=0
x=x+2
x=x+1
x=x+1
x=x+2
x=0
x=0
x=x+3
x=0
考虑数据
x=0
x=0
x=x+2
x=x+3 | Legal
x=x+3 | Legal
x=x+3 | Illegal
x=x+1
x=x+2 | Illegal
\end{verbatim}

\subsection{最终一致性}

最终一致性 (Eventual Consistency) 定义：
在某一时刻停止更新后，所有更新都会传播，使得所有副本存储相同数据 (直到再次出现更新)。

直观例子：分布式系统存储集合和并发写操作。最终所有副本收敛一致。

一致性模型对比：

\begin{table}[htbp]
\centering
\caption{一致性模型对比}
\begin{tabular}{llll}
\toprule
一致性模型 & 核心特征 & 示例 & 强度 \\
\midrule
线性一致性 & 操作实时性约束，符合真实时间，严格原子生效，跨进程 & 银行转账: ATM取款后，全球查询余额立刻同步。 & 极强 (Real-time) \\
顺序一致性 & 所有节点看到相同的全局总序，且不违反进程内顺序。 & 网络游戏：所有玩家看到的击杀顺序必须一致，虽然可能有延迟。 & 中 (Logical-time) \\
因果一致性 & 仅保证有因果关系的写操作在所有节点顺序一致。 & 论坛发帖：先看到主贴 (因)，后看到回复 (果)。 & 中 (Causal-order) \\
最终一致性 & 不保证中间读取得结果，但系统最终会收敛到一致。 & DNS 修改：全球各处生效时间不一。 & 弱 (No guarantee) \\
\bottomrule
\end{tabular}
\end{table}

\section{客户端一致性模型}

核心思想：不追求全局一致，只保证单个客户端访问这个分布式数据库时，在任意位置读取得结果与之前在另一位置离开时的状态一致，即数据对该客户端而言是一致的。具体包括以下内容：
\begin{itemize}
\item 单调读：读不会“时光倒流” (应当越读越新)
\item 单调写：写不会“顺序颠倒” (应当先写先到)
\item 读你所写：自己写的，自己立刻能读到
\item 读后写：基于读的写，存在于读之后更新读版本 (因果)
\end{itemize}

客户端一致性实现基本思想：
\begin{itemize}
\item 每个写操作由其源服务器分配全局唯一标识符
\item 为每个客户端跟踪两个写操作集：读集和写集。
\end{itemize}

实现：
\begin{itemize}
\item 单调读：客户端传递读集，服务器拉取更新并更新读集。
\item 单调写：客户端传递写集，服务器拉取更新、排序并执行写操作，更新写集。
\item 读你所写：客户端传递写集，服务器拉取更新并执行读操作，更新读集。
\item 读后写：客户端传递读集，服务器拉取更新、排序并执行写操作，更新写集。
\end{itemize}

\section{副本管理 (对应 7.4 Replica management)}

\subsection{副本放置 (Replica Placement)}

目标：从 $N$ 个可能位置中找出最佳的 $K$ 个位置。
策略：
\begin{itemize}
\item 选择平均客户端距离最小的位置。(计算量大)
\item 选择第 $K$ 大自治系统，并在连接最好的主机上放置服务器。(计算量大)
\item 在反映延迟的 $d$ 维几何空间中定位节点，识别密度最高的 $K$ 个区域并在每个区域放置服务器。(计算量小)
\end{itemize}

\subsection{内容复制 (Content Replication)}

副本类型：
\begin{itemize}
\item 永久副本 (Permanent replicas)：始终拥有副本的进程/机器，可以被视为构成分布式数据存储的初始副本集。
\item 服务器发起副本 (Server-initiated replica)：是数据存储的副本，其存在是为了提升性能，并且是由数据存储的所有者主动创建的。托管副本设施，客户端使用它。
\item 客户端发起副本 (Client-initiated replica)：根据客户端请求动态的进程 (客户端缓存)。本质上，缓存是一种本地存储来临时存储刚刚请求的数据副本。
\end{itemize}

\subsection{内容分发 (Content Distribution)}

方式：
\begin{itemize}
\item 传播更新通知/失效 (常用于缓存)。
\item 在副本间传输数据 (分布式数据库)：被动复制。
\item 传播更新操作到其他副本：主动复制。
\end{itemize}

推送 (Push) 与拉取 (Pull)：
\begin{itemize}
\item 推送：服务器发起，无论目标是否请求。
\item 拉取：客户端发起，客户端请求更新。
\end{itemize}

租约 (Leases)：
\begin{itemize}
\item 租约是一种契约，服务器承诺在租约到期前将更新推送给客户端。
\item 自适应租约到期时间：基于年龄、续订频率、服务器状态。
\end{itemize}

\subsection{管理复制对象}

目标：
\begin{itemize}
\item 防止对同一对象的多次调用并发执行。
\item 确保复制状态的所有更改都相同。
\end{itemize}

解决方案：
\begin{itemize}
\item 使用本地锁定机制序列化对象内部数据访问。
\item 需要确定性线程调度以确保状态更改一致。
\end{itemize}

汇报完毕，欢迎同学们批评指教！本次课大致对应课程教材第七章。

% 原文可能存在不完整语句：“库对你而言是一致的。具有节 | [疑似: 未完]”，已忽略。
